\documentclass[11pt,a4paper]{article}
\usepackage[T1]{fontenc}
\usepackage[utf8]{inputenc}
\usepackage{amsmath,amssymb,xcolor}
\usepackage[margin=25mm]{geometry}
\usepackage[colorlinks=true,urlcolor=blue]{hyperref}
\definecolor{nexorablue}{HTML}{2459B5}
\setlength{\parindent}{0pt}
\setlength{\parskip}{7pt}
\setlength{\emergencystretch}{2em}
\renewcommand{\familydefault}{\sfdefault}
\begin{document}

{\small\bfseries\color{nexorablue} NEXORA CPGE}\par
{\small MP / NEX-STR-03 / Corrigé}\par
{\LARGE\bfseries Une congruence à coefficient non inversible}\par
\vspace{0.25em}\hrule\vspace{0.6em}

\section*{Module réduit}
L’égalité équivaut à $2x-3=5k$. L’inverse de $2$ modulo $5$ vaut $3$, donc $x\equiv4\pmod5$.

\section*{Retour au module initial}
Les classes modulo $35$ sont représentées par $4,9,14,19,24,29,34$. Avec $22$, l’équation est impossible, car les multiples de $14$ et de $35$ sont divisibles par $7$, contrairement à $22$.


\end{document}
